as in Fulton--Harris \cite{FH} or Serre~\cite{serre}.
We will make use of a standard character-theoretic technique based on the following result due to Frobenius.
For a straightforward proof, see Zagier's Appendix A in Lando--Zvonkin \cite{LZ}.
Let \(G\) be a finite group, and let \(\Irr (G)\) denote the set of all irreducible characters of \(G\).
For \(\chi \in \Irr (G)\), let \(\deg \chi\) denote the value of \(\chi\) at the identity element of \(G\).


\begin{theorem}[Frobenius \cite{frob}]\label{frob}
Let \(k\) be a positive integer, and, for each \(i \in \{1,\ldots,k\}\), let \(A_i\) be a union of conjugacy classes in \(G\).
For any \(g \in G\), the number of tuples \((t_1,\ldots,t_k) \in A_1 \times \cdots \times A_k\) such that \(t_1 \cdots t_k = g\) is given by
\begin{equation}
\frac{1}{\# G} \sum_{\chi \in \Irr (G)} (\deg \chi)^{1-k} \chi (g^{-1}) \prod_{i = 1}^k \sum_{t \in A_i} \chi(t).
\label{ff}
\end{equation}
\end{theorem}


\begin{corollary}
\label{ff_in_situ_corollary}
For all \(n,k \in \N\), \(\mu \vdash n\), and prime powers~\(q\),
\begin{equation}
g_{k,\mu} (q)
=
\label{ff_in_situ}
\frac{1}{\# \GL_n \F_q}
\sum_{\chi \in \Irr (\GL_n \F_q)}
(\deg \chi)^{1-k}
\left ( \sum_{g \in \CT_{(n)} (q)} \chi(g) \right )^k
\left ( \sum_{h \in \CT_\mu (q)}   \chi(h) \right ).
\end{equation}
